\begin{frame}{확인 문제}
  \begin{enumerate}
    \item \kb{공유 테이블.} \texttt{dyna\_q\_grid}에서 계획 갱신(3번
      단계)을 \texttt{Q\_plan} 별도 딕셔너리로 분리하면,
      \texttt{planning\_steps=30}의 수렴 에피소드가 순수
      Q-learning(334)과 어떻게 달라지는지 \textbf{숫자 없이}
      한 문장. (힌트: 계획의 산출물이 어디로 가는가.)
    \item \kb{모델 커버리지.} 400 에피소드 후에도 100개 $(s,a)$ 중
      4개(4\%)가 모델에 안 들어감. 이 ``없는 4개''가 계획에
      미치는 영향을 \texttt{random.choice(list(model.items()))}의
      관점에서 설명.
    \item \kb{참값 판정.} 학습 후 $Q$에서
      $\max(Q[(\text{START}, 0\sim3)]) = -7.00$이고
      $\arg\max$가 ``위''(action 0) 또는 ``오른쪽''(action 3)이면,
      이 정책이 \textbf{최적}인 근거가 손계산 참값(위/오른쪽 $=-7$,
      아래/왼쪽 $=-8$)에 어떻게 대응하는지 설명.
  \end{enumerate}
\end{frame}
